\SourceLeaf{293}%
the following rearrangement is most advantageous for a direct, rigorous solution. The first two formulas remain the same as in the first method.

Fourth method
\[
\begin{aligned}
\operatorname{tang}s&=\cos T\operatorname{tang}r\\
\operatorname{tang}\lambda&=\frac{\operatorname{tang}T\sin s}{\cos(S-s)}\\
\operatorname{tang}t&=\sin T\sin r\operatorname{tang}(S-s)\\
\sin\tau&=\sin T\operatorname{tang}\tfrac{1}{2}r\sin s\\
\sin\sigma&=\operatorname{tang}t\operatorname{tang}\tfrac{1}{2}\lambda\cos(S-s)\\
S'&=S-s-\sigma\\
T'&=T-t-\tau
\end{aligned}
\]
\indent This fourth method leaves nothing to be desired in accuracy. Calculating directly in this form, however, requires rather laborious interpolation to determine small arcs from the logarithms of their tangents or sines. This drawback is easily avoided by expanding the trigonometric functions in series. We can then carry out the calculations using the logarithms of numbers without losing accuracy. It will suffice to give only the main steps of this transformation here.

Let
\[
\begin{aligned}
r\cos T&=s^0\\
r\sin T&=v
\end{aligned}
\]
Then, if for brevity we denote the length of one second of arc in units of the radius, or the fraction \(\tfrac{\pi}{648000}\), by \(\rho\) and regard \(r\) as a quantity of first order, we have, with accuracy up to but excluding quantities of fifth order,
\[
s=s^0(1+\tfrac{1}{3}\rho\rho rr-\tfrac{1}{3}\rho\rho s^0s^0)=s^0(1+\tfrac{1}{3}\rho\rho vv)
\]
If we then also set
\[
\begin{aligned}
v\operatorname{tang}(S-s)&=t^0\\
\frac{v}{\cos(S-s)}&=\lambda^0
\end{aligned}
\]
we obtain
